% --- Archon physics formalization source begin ---
% archon:physics
% archon:covers PhyXMiniProblems/problem_phyx_mini_0529.lean
% archon:source-report reports/phyx_mini/problem_phyx_mini_0529.source.json

\chapter{Physics problem phyx\_mini\_0529}
\label{ch:PhyXMiniProblems_problem_phyx_mini_0529}

\paragraph{Problem source.}
\#\# Physical scenario

Owen and Dina are at rest in frame S\textbackslash{}( '\textbackslash{}), which is moving at \textbackslash{}( 0.600c \textbackslash{}) with respect to frame S. They play a game of catch while Ed, at rest in frame S, watches the action as shown in figure. Owen throws the ball to Dina at \textbackslash{}( 0.800c \textbackslash{}) (according to Owen), and their separation (measured in S\textbackslash{}( '\textbackslash{})) is equal to \textbackslash{}( 1.80 \textbackslash{}times 10\textasciicircum{}\{12\} \textbackslash{}) m.

\#\# Question

what time interval is required for the ball to reach Dina?

\#\# Answer choices

- A: A: \textbackslash{}( 3.22 \textbackslash{}times 10\textasciicircum{}3 \textbackslash{}text\{ s\} \textbackslash{})

- B: B: \textbackslash{}( 2.54 \textbackslash{}times 10\textasciicircum{}3 \textbackslash{}text\{ s\} \textbackslash{})

- C: C: \textbackslash{}( 3.71 \textbackslash{}times 10\textasciicircum{}3 \textbackslash{}text\{ s\} \textbackslash{})

- D: D: \textbackslash{}( 4.88 \textbackslash{}times 10\textasciicircum{}3 \textbackslash{}text\{ s\} \textbackslash{})

\#\# Figure caption (auxiliary; use the image as primary evidence)

The image presents a diagram with three observers: Dina, Owen, and Ed. It depicts a relativistic scenario with the following details:

1. **Observers and Positions**:
   - Dina is located in the \textbackslash{}(S'\textbackslash{}) reference frame, with an arrow indicating a velocity of \textbackslash{}(0.600c\textbackslash{}) directed to the right.
   - Owen is positioned towards the right of Dina in the \textbackslash{}(x'\textbackslash{}) reference frame, appearing to throw a ball to the left.
   - Ed is below Dina and Owen, located in the ground frame \textbackslash{}(S\textbackslash{}).

2. **Ball**:
   - Owen throws a ball, moving to the left with a velocity of \textbackslash{}(0.800c\textbackslash{}).

3. **Distances**:
   - A line along the \textbackslash{}(x'\textbackslash{}) axis marks the distance of \textbackslash{}(1.80 \textbackslash{}times 10\textasciicircum{}\{12\}\textbackslash{}) meters between Dina and some point further to the right.

4. **Axes**:
   - The diagram includes two coordinate axes, \textbackslash{}(S\textbackslash{}) and \textbackslash{}(x\textbackslash{}), forming a traditional Cartesian plane with \textbackslash{}(x\textbackslash{}) running horizontally. Dina and Owen are shown related to a dashed line indicating movement in the \textbackslash{}(S'\textbackslash{}) frame.

5. **Vectors**:
   - Red arrows represent velocities for both Dina and the ball, indicating their directions and speeds.

This illustration likely relates to concepts in special relativity, demonstrating relative velocities between observers in different frames of reference.

\#\# Dataset metadata

Relativity; Physical Model Grounding Reasoning, Multi-Formula Reasoning

\paragraph{Recorded answer/context.}
D: D: \textbackslash{}( 4.88 \textbackslash{}times 10\textasciicircum{}3 \textbackslash{}text\{ s\} \textbackslash{})

\paragraph{Figure/image path.}
/root/proposal\_for\_physic/science-mango/phyx\_mini\_run/phyx\_data/test\_image/529.png

\paragraph{Formalization target.}
create a compiling Lean file with sorry bodies at `PhyXMiniProblems/problem\_phyx\_mini\_0529.lean`. The Lean declarations must preserve the physical quantities, dimensions or dimensional roles, figure labels, governing-law hypotheses, and final relation expressed by this problem.
Use Mathlib/Physlib names found through LeanExplore where available. If a physics API is missing, introduce faithful local abstractions rather than scalar placeholder aliases.

\begin{theorem}[Physics formalization target]
\label{thm:physics:phyx_mini_0529:target}
\lean{PhyXMiniProblems.ProblemPhyXMini0529.ballFlightTimeInEdFrame_matches_answerD}
\uses{def:physics:phyx-mini-0529:phyxminiproblems-problemphyxmini0529-relativisticcatchsetup, def:physics:phyx-mini-0529:phyxminiproblems-problemphyxmini0529-matchesscenarioandprimaryfigure, def:physics:phyx-mini-0529:phyxminiproblems-problemphyxmini0529-hasphysicalrelativisticparameters, def:physics:phyx-mini-0529:phyxminiproblems-problemphyxmini0529-satisfiesuniformballmotioninmovingframe, def:physics:phyx-mini-0529:phyxminiproblems-problemphyxmini0529-satisfieslorentztimetransformation, def:physics:phyx-mini-0529:phyxminiproblems-problemphyxmini0529-edmeasuresgroundframeelapsedtime, def:physics:phyx-mini-0529:phyxminiproblems-problemphyxmini0529-recordeddatasetanswer, def:physics:phyx-mini-0529:phyxminiproblems-problemphyxmini0529-matchesdisplayedtimechoice}
The assigned autoformalize agent should translate this physics problem into Lean declarations in the covered file, with theorem and lemma proof bodies written as `by sorry`.
\end{theorem}
\begin{proof}
This is an autoformalization task, not a proof task. Produce statements that can later be proved by the physics prover without weakening the physical model.
\end{proof}

% --- Archon declaration topology begin ---
\section{Declaration topology}
\begin{definition}[Length Quantity]
  \label{def:physics:phyx-mini-0529:phyxminiproblems-problemphyxmini0529-lengthquantity}
  \lean{PhyXMiniProblems.ProblemPhyXMini0529.LengthQuantity}
  A nonnegative physical length, independent of the selected unit
\end{definition}
\begin{proof}
  This is the declared model definition of Length Quantity.
\end{proof}
\begin{definition}[Duration Quantity]
  \label{def:physics:phyx-mini-0529:phyxminiproblems-problemphyxmini0529-durationquantity}
  \lean{PhyXMiniProblems.ProblemPhyXMini0529.DurationQuantity}
  A nonnegative physical elapsed time
\end{definition}
\begin{proof}
  This is the declared model definition of Duration Quantity.
\end{proof}
\begin{definition}[Axial Coordinate]
  \label{def:physics:phyx-mini-0529:phyxminiproblems-problemphyxmini0529-axialcoordinate}
  \lean{PhyXMiniProblems.ProblemPhyXMini0529.AxialCoordinate}
  A signed coordinate along the common one-dimensional spatial axis
\end{definition}
\begin{proof}
  This is the declared model definition of Axial Coordinate.
\end{proof}
\begin{definition}[Clock Coordinate]
  \label{def:physics:phyx-mini-0529:phyxminiproblems-problemphyxmini0529-clockcoordinate}
  \lean{PhyXMiniProblems.ProblemPhyXMini0529.ClockCoordinate}
  A signed clock-coordinate value in an inertial frame
\end{definition}
\begin{proof}
  This is the declared model definition of Clock Coordinate.
\end{proof}
\begin{definition}[Axial Velocity]
  \label{def:physics:phyx-mini-0529:phyxminiproblems-problemphyxmini0529-axialvelocity}
  \lean{PhyXMiniProblems.ProblemPhyXMini0529.AxialVelocity}
  A signed physical velocity along the common spatial axis
\end{definition}
\begin{proof}
  This is the declared model definition of Axial Velocity.
\end{proof}
\begin{definition}[length Readout]
  \label{def:physics:phyx-mini-0529:phyxminiproblems-problemphyxmini0529-lengthreadout}
  \lean{PhyXMiniProblems.ProblemPhyXMini0529.lengthReadout}
  \uses{def:physics:phyx-mini-0529:phyxminiproblems-problemphyxmini0529-lengthquantity}
  Read a nonnegative physical length in a selected length unit
\end{definition}
\begin{proof}
  This is the declared model definition of length Readout.
\end{proof}
\begin{definition}[duration Readout]
  \label{def:physics:phyx-mini-0529:phyxminiproblems-problemphyxmini0529-durationreadout}
  \lean{PhyXMiniProblems.ProblemPhyXMini0529.durationReadout}
  \uses{def:physics:phyx-mini-0529:phyxminiproblems-problemphyxmini0529-durationquantity}
  Read a nonnegative physical duration in a selected time unit
\end{definition}
\begin{proof}
  This is the declared model definition of duration Readout.
\end{proof}
\begin{definition}[axial Coordinate Readout]
  \label{def:physics:phyx-mini-0529:phyxminiproblems-problemphyxmini0529-axialcoordinatereadout}
  \lean{PhyXMiniProblems.ProblemPhyXMini0529.axialCoordinateReadout}
  \uses{def:physics:phyx-mini-0529:phyxminiproblems-problemphyxmini0529-axialcoordinate}
  Read a signed axial coordinate in a selected length unit
\end{definition}
\begin{proof}
  This is the declared model definition of axial Coordinate Readout.
\end{proof}
\begin{definition}[clock Coordinate Readout]
  \label{def:physics:phyx-mini-0529:phyxminiproblems-problemphyxmini0529-clockcoordinatereadout}
  \lean{PhyXMiniProblems.ProblemPhyXMini0529.clockCoordinateReadout}
  \uses{def:physics:phyx-mini-0529:phyxminiproblems-problemphyxmini0529-clockcoordinate}
  Read a signed clock coordinate in a selected time unit
\end{definition}
\begin{proof}
  This is the declared model definition of clock Coordinate Readout.
\end{proof}
\begin{definition}[axial Velocity Readout]
  \label{def:physics:phyx-mini-0529:phyxminiproblems-problemphyxmini0529-axialvelocityreadout}
  \lean{PhyXMiniProblems.ProblemPhyXMini0529.axialVelocityReadout}
  \uses{def:physics:phyx-mini-0529:phyxminiproblems-problemphyxmini0529-axialvelocity}
  Read a signed axial velocity in coherent selected length and time units
\end{definition}
\begin{proof}
  This is the declared model definition of axial Velocity Readout.
\end{proof}
\begin{definition}[length In Meters]
  \label{def:physics:phyx-mini-0529:phyxminiproblems-problemphyxmini0529-lengthinmeters}
  \lean{PhyXMiniProblems.ProblemPhyXMini0529.lengthInMeters}
  \uses{def:physics:phyx-mini-0529:phyxminiproblems-problemphyxmini0529-lengthquantity, def:physics:phyx-mini-0529:phyxminiproblems-problemphyxmini0529-lengthreadout}
  Metre readout used by the separation and spatial event coordinates
\end{definition}
\begin{proof}
  This is the declared model definition of length In Meters.
\end{proof}
\begin{definition}[axial Coordinate In Meters]
  \label{def:physics:phyx-mini-0529:phyxminiproblems-problemphyxmini0529-axialcoordinateinmeters}
  \lean{PhyXMiniProblems.ProblemPhyXMini0529.axialCoordinateInMeters}
  \uses{def:physics:phyx-mini-0529:phyxminiproblems-problemphyxmini0529-axialcoordinate, def:physics:phyx-mini-0529:phyxminiproblems-problemphyxmini0529-axialcoordinatereadout}
  Metre readout of a signed axial coordinate
\end{definition}
\begin{proof}
  This is the declared model definition of axial Coordinate In Meters.
\end{proof}
\begin{definition}[duration In Seconds]
  \label{def:physics:phyx-mini-0529:phyxminiproblems-problemphyxmini0529-durationinseconds}
  \lean{PhyXMiniProblems.ProblemPhyXMini0529.durationInSeconds}
  \uses{def:physics:phyx-mini-0529:phyxminiproblems-problemphyxmini0529-durationquantity, def:physics:phyx-mini-0529:phyxminiproblems-problemphyxmini0529-durationreadout}
  Second readout of a physical elapsed time
\end{definition}
\begin{proof}
  This is the declared model definition of duration In Seconds.
\end{proof}
\begin{definition}[clock Coordinate In Seconds]
  \label{def:physics:phyx-mini-0529:phyxminiproblems-problemphyxmini0529-clockcoordinateinseconds}
  \lean{PhyXMiniProblems.ProblemPhyXMini0529.clockCoordinateInSeconds}
  \uses{def:physics:phyx-mini-0529:phyxminiproblems-problemphyxmini0529-clockcoordinate, def:physics:phyx-mini-0529:phyxminiproblems-problemphyxmini0529-clockcoordinatereadout}
  Second readout of a signed inertial-frame clock coordinate
\end{definition}
\begin{proof}
  This is the declared model definition of clock Coordinate In Seconds.
\end{proof}
\begin{definition}[axial Velocity In Meters Per Second]
  \label{def:physics:phyx-mini-0529:phyxminiproblems-problemphyxmini0529-axialvelocityinmeterspersecond}
  \lean{PhyXMiniProblems.ProblemPhyXMini0529.axialVelocityInMetersPerSecond}
  \uses{def:physics:phyx-mini-0529:phyxminiproblems-problemphyxmini0529-axialvelocity, def:physics:phyx-mini-0529:phyxminiproblems-problemphyxmini0529-axialvelocityreadout}
  Metres-per-second readout of a signed axial velocity
\end{definition}
\begin{proof}
  This is the declared model definition of axial Velocity In Meters Per Second.
\end{proof}
\begin{definition}[vacuum Speed Of Light Readout]
  \label{def:physics:phyx-mini-0529:phyxminiproblems-problemphyxmini0529-vacuumspeedoflightreadout}
  \lean{PhyXMiniProblems.ProblemPhyXMini0529.vacuumSpeedOfLightReadout}
  Physlib's exact vacuum speed of light in selected length and time units
\end{definition}
\begin{proof}
  This is the declared model definition of vacuum Speed Of Light Readout.
\end{proof}
\begin{definition}[velocity In Light Speed Units]
  \label{def:physics:phyx-mini-0529:phyxminiproblems-problemphyxmini0529-velocityinlightspeedunits}
  \lean{PhyXMiniProblems.ProblemPhyXMini0529.velocityInLightSpeedUnits}
  \uses{def:physics:phyx-mini-0529:phyxminiproblems-problemphyxmini0529-axialvelocity, def:physics:phyx-mini-0529:phyxminiproblems-problemphyxmini0529-axialvelocityinmeterspersecond, def:physics:phyx-mini-0529:phyxminiproblems-problemphyxmini0529-vacuumspeedoflightreadout}
  The signed dimensionless ratio v / c, evaluated in coherent SI units
\end{definition}
\begin{proof}
  This is the declared model definition of velocity In Light Speed Units.
\end{proof}
\begin{definition}[Inertial Frame Label]
  \label{def:physics:phyx-mini-0529:phyxminiproblems-problemphyxmini0529-inertialframelabel}
  \lean{PhyXMiniProblems.ProblemPhyXMini0529.InertialFrameLabel}
  The two inertial frames named in the problem and figure
\end{definition}
\begin{proof}
  This is the declared model definition of Inertial Frame Label.
\end{proof}
\begin{definition}[Coordinate Axis Label]
  \label{def:physics:phyx-mini-0529:phyxminiproblems-problemphyxmini0529-coordinateaxislabel}
  \lean{PhyXMiniProblems.ProblemPhyXMini0529.CoordinateAxisLabel}
  The horizontal coordinate axes printed for the two frames
\end{definition}
\begin{proof}
  This is the declared model definition of Coordinate Axis Label.
\end{proof}
\begin{definition}[Observer Label]
  \label{def:physics:phyx-mini-0529:phyxminiproblems-problemphyxmini0529-observerlabel}
  \lean{PhyXMiniProblems.ProblemPhyXMini0529.ObserverLabel}
  The three observers named in the scenario and shown in the figure
\end{definition}
\begin{proof}
  This is the declared model definition of Observer Label.
\end{proof}
\begin{definition}[Catch Event]
  \label{def:physics:phyx-mini-0529:phyxminiproblems-problemphyxmini0529-catchevent}
  \lean{PhyXMiniProblems.ProblemPhyXMini0529.CatchEvent}
  The two endpoint events of the ball's flight
\end{definition}
\begin{proof}
  This is the declared model definition of Catch Event.
\end{proof}
\begin{definition}[Axial Direction]
  \label{def:physics:phyx-mini-0529:phyxminiproblems-problemphyxmini0529-axialdirection}
  \lean{PhyXMiniProblems.ProblemPhyXMini0529.AxialDirection}
  Orientations along the common horizontal axis
\end{definition}
\begin{proof}
  This is the declared model definition of Axial Direction.
\end{proof}
\begin{definition}[Figure Feature]
  \label{def:physics:phyx-mini-0529:phyxminiproblems-problemphyxmini0529-figurefeature}
  \lean{PhyXMiniProblems.ProblemPhyXMini0529.FigureFeature}
  Individually identifiable items present in the supplied primary image
\end{definition}
\begin{proof}
  This is the declared model definition of Figure Feature.
\end{proof}
\begin{definition}[Relativistic Catch Figure]
  \label{def:physics:phyx-mini-0529:phyxminiproblems-problemphyxmini0529-relativisticcatchfigure}
  \lean{PhyXMiniProblems.ProblemPhyXMini0529.RelativisticCatchFigure}
  \uses{def:physics:phyx-mini-0529:phyxminiproblems-problemphyxmini0529-inertialframelabel, def:physics:phyx-mini-0529:phyxminiproblems-problemphyxmini0529-coordinateaxislabel, def:physics:phyx-mini-0529:phyxminiproblems-problemphyxmini0529-observerlabel, def:physics:phyx-mini-0529:phyxminiproblems-problemphyxmini0529-axialdirection, def:physics:phyx-mini-0529:phyxminiproblems-problemphyxmini0529-figurefeature}
  Qualitative geometry and scalar annotations read directly from phyx\_data/test\_image/529.png
\end{definition}
\begin{proof}
  This is the declared model definition of Relativistic Catch Figure.
\end{proof}
\begin{definition}[Relativistic Catch Setup]
  \label{def:physics:phyx-mini-0529:phyxminiproblems-problemphyxmini0529-relativisticcatchsetup}
  \lean{PhyXMiniProblems.ProblemPhyXMini0529.RelativisticCatchSetup}
  \uses{def:physics:phyx-mini-0529:phyxminiproblems-problemphyxmini0529-lengthquantity, def:physics:phyx-mini-0529:phyxminiproblems-problemphyxmini0529-durationquantity, def:physics:phyx-mini-0529:phyxminiproblems-problemphyxmini0529-axialcoordinate, def:physics:phyx-mini-0529:phyxminiproblems-problemphyxmini0529-clockcoordinate, def:physics:phyx-mini-0529:phyxminiproblems-problemphyxmini0529-axialvelocity, def:physics:phyx-mini-0529:phyxminiproblems-problemphyxmini0529-inertialframelabel, def:physics:phyx-mini-0529:phyxminiproblems-problemphyxmini0529-observerlabel, def:physics:phyx-mini-0529:phyxminiproblems-problemphyxmini0529-catchevent, def:physics:phyx-mini-0529:phyxminiproblems-problemphyxmini0529-relativisticcatchfigure}
  All independent physical quantities and frame-indexed event coordinates in the scenario. edMeasuredFlightTime is an unknown duration; no field assigns it a displayed answer value
\end{definition}
\begin{proof}
  This is the declared model definition of Relativistic Catch Setup.
\end{proof}
\begin{definition}[Matches Scenario And Primary Figure]
  \label{def:physics:phyx-mini-0529:phyxminiproblems-problemphyxmini0529-matchesscenarioandprimaryfigure}
  \lean{PhyXMiniProblems.ProblemPhyXMini0529.MatchesScenarioAndPrimaryFigure}
  \uses{def:physics:phyx-mini-0529:phyxminiproblems-problemphyxmini0529-lengthinmeters, def:physics:phyx-mini-0529:phyxminiproblems-problemphyxmini0529-axialcoordinateinmeters, def:physics:phyx-mini-0529:phyxminiproblems-problemphyxmini0529-velocityinlightspeedunits, def:physics:phyx-mini-0529:phyxminiproblems-problemphyxmini0529-relativisticcatchsetup}
  The named labels, directions, and numerical annotations in the source and primary image. Positive axial velocity points right, so the ball's signed S' velocity is the negative of its printed speed magnitude. The two event positions express that the throw occurs at Owen and the catch at Dina. Their difference is fixed by the independently represented proper separation; no elapsed-time answer is recorded here
\end{definition}
\begin{proof}
  This is the declared model definition of Matches Scenario And Primary Figure.
\end{proof}
\begin{definition}[Has Physical Relativistic Parameters]
  \label{def:physics:phyx-mini-0529:phyxminiproblems-problemphyxmini0529-hasphysicalrelativisticparameters}
  \lean{PhyXMiniProblems.ProblemPhyXMini0529.HasPhysicalRelativisticParameters}
  \uses{def:physics:phyx-mini-0529:phyxminiproblems-problemphyxmini0529-lengthinmeters, def:physics:phyx-mini-0529:phyxminiproblems-problemphyxmini0529-durationinseconds, def:physics:phyx-mini-0529:phyxminiproblems-problemphyxmini0529-clockcoordinateinseconds, def:physics:phyx-mini-0529:phyxminiproblems-problemphyxmini0529-vacuumspeedoflightreadout, def:physics:phyx-mini-0529:phyxminiproblems-problemphyxmini0529-velocityinlightspeedunits, def:physics:phyx-mini-0529:phyxminiproblems-problemphyxmini0529-relativisticcatchsetup}
  Positivity, causal speed bounds, and event ordering for the physical setup. These conditions contain no displayed flight-time value
\end{definition}
\begin{proof}
  This is the declared model definition of Has Physical Relativistic Parameters.
\end{proof}
\begin{definition}[Satisfies Uniform Ball Motion In Moving Frame]
  \label{def:physics:phyx-mini-0529:phyxminiproblems-problemphyxmini0529-satisfiesuniformballmotioninmovingframe}
  \lean{PhyXMiniProblems.ProblemPhyXMini0529.SatisfiesUniformBallMotionInMovingFrame}
  \uses{def:physics:phyx-mini-0529:phyxminiproblems-problemphyxmini0529-axialcoordinateinmeters, def:physics:phyx-mini-0529:phyxminiproblems-problemphyxmini0529-clockcoordinateinseconds, def:physics:phyx-mini-0529:phyxminiproblems-problemphyxmini0529-axialvelocityinmeterspersecond, def:physics:phyx-mini-0529:phyxminiproblems-problemphyxmini0529-relativisticcatchsetup}
  Uniform one-dimensional motion of the ball in Owen and Dina's rest frame: Delta x' = u' Delta t'. This is a generic kinematic law involving the unknown event-time difference, not the requested numerical answer
\end{definition}
\begin{proof}
  This is the declared model definition of Satisfies Uniform Ball Motion In Moving Frame.
\end{proof}
\begin{definition}[Satisfies Lorentz Time Transformation]
  \label{def:physics:phyx-mini-0529:phyxminiproblems-problemphyxmini0529-satisfieslorentztimetransformation}
  \lean{PhyXMiniProblems.ProblemPhyXMini0529.SatisfiesLorentzTimeTransformation}
  \uses{def:physics:phyx-mini-0529:phyxminiproblems-problemphyxmini0529-axialcoordinateinmeters, def:physics:phyx-mini-0529:phyxminiproblems-problemphyxmini0529-clockcoordinateinseconds, def:physics:phyx-mini-0529:phyxminiproblems-problemphyxmini0529-vacuumspeedoflightreadout, def:physics:phyx-mini-0529:phyxminiproblems-problemphyxmini0529-velocityinlightspeedunits, def:physics:phyx-mini-0529:phyxminiproblems-problemphyxmini0529-relativisticcatchsetup}
  The longitudinal Lorentz transformation for the separation of the catch and throw events, Delta t = gamma(beta) * (Delta t' + beta * Delta x' / c). Here beta is the signed velocity of S' in S; the leftward event displacement Delta x' supplies the minus sign relevant to this figure. The law is stated before inserting or rounding any answer value
\end{definition}
\begin{proof}
  This is the declared model definition of Satisfies Lorentz Time Transformation.
\end{proof}
\begin{definition}[Ed Measures Ground Frame Elapsed Time]
  \label{def:physics:phyx-mini-0529:phyxminiproblems-problemphyxmini0529-edmeasuresgroundframeelapsedtime}
  \lean{PhyXMiniProblems.ProblemPhyXMini0529.EdMeasuresGroundFrameElapsedTime}
  \uses{def:physics:phyx-mini-0529:phyxminiproblems-problemphyxmini0529-durationinseconds, def:physics:phyx-mini-0529:phyxminiproblems-problemphyxmini0529-clockcoordinateinseconds, def:physics:phyx-mini-0529:phyxminiproblems-problemphyxmini0529-relativisticcatchsetup}
  Operational meaning of the question: because Ed is at rest in S, his measured flight duration is the ground-frame coordinate-time difference. This relation does not assign that duration a numerical answer
\end{definition}
\begin{proof}
  This is the declared model definition of Ed Measures Ground Frame Elapsed Time.
\end{proof}
\begin{definition}[Answer Choice]
  \label{def:physics:phyx-mini-0529:phyxminiproblems-problemphyxmini0529-answerchoice}
  \lean{PhyXMiniProblems.ProblemPhyXMini0529.AnswerChoice}
  The four answer labels printed with the problem
\end{definition}
\begin{proof}
  This is the declared model definition of Answer Choice.
\end{proof}
\begin{definition}[displayed Time Seconds]
  \label{def:physics:phyx-mini-0529:phyxminiproblems-problemphyxmini0529-displayedtimeseconds}
  \lean{PhyXMiniProblems.ProblemPhyXMini0529.displayedTimeSeconds}
  \uses{def:physics:phyx-mini-0529:phyxminiproblems-problemphyxmini0529-answerchoice}
  Time in seconds printed beside each multiple-choice label
\end{definition}
\begin{proof}
  This is the declared model definition of displayed Time Seconds.
\end{proof}
\begin{definition}[recorded Dataset Answer]
  \label{def:physics:phyx-mini-0529:phyxminiproblems-problemphyxmini0529-recordeddatasetanswer}
  \lean{PhyXMiniProblems.ProblemPhyXMini0529.recordedDatasetAnswer}
  \uses{def:physics:phyx-mini-0529:phyxminiproblems-problemphyxmini0529-answerchoice}
  The answer label recorded in the supplied dataset metadata
\end{definition}
\begin{proof}
  This is the declared model definition of recorded Dataset Answer.
\end{proof}
\begin{definition}[Matches Displayed Time Choice]
  \label{def:physics:phyx-mini-0529:phyxminiproblems-problemphyxmini0529-matchesdisplayedtimechoice}
  \lean{PhyXMiniProblems.ProblemPhyXMini0529.MatchesDisplayedTimeChoice}
  \uses{def:physics:phyx-mini-0529:phyxminiproblems-problemphyxmini0529-durationquantity, def:physics:phyx-mini-0529:phyxminiproblems-problemphyxmini0529-durationinseconds, def:physics:phyx-mini-0529:phyxminiproblems-problemphyxmini0529-answerchoice, def:physics:phyx-mini-0529:phyxminiproblems-problemphyxmini0529-displayedtimeseconds}
  Agreement with a time printed as three significant figures in units of 10\textasciicircum{}3 s. Half of the final 10 s place gives a 5 s rounding tolerance
\end{definition}
\begin{proof}
  This is the declared model definition of Matches Displayed Time Choice.
\end{proof}
% --- Archon declaration topology end ---
% --- Archon physics formalization source end ---
